%! TeX program = lualatex
% Converted from book/index.md by book/tools/md2tex.py (initial conversion).
% This .tex file is the maintained source from now on; the converter is kept for provenance only.
\documentclass[../book.tex]{subfiles}
\begin{document}

\chapter*{Welcome to Applied Math 1201 @ UWO}\label{ch:welcome}
\addcontentsline{toc}{chapter}{Welcome to Applied Math 1201 @ UWO}

\begin{attentionbox}[Read the course outline on \href{https://owl.uwo.ca}{OWL} first!]\label{warn:owl}
This website publishes the course textbook and practice problems.

If this is your \emph{first} visit here and \hlmain{have not read the course outline}, please stop and go to \href{https://owl.uwo.ca}{OWL} for the course outline, assignments and additional course material.
\end{attentionbox}
Before you proceed, please make sure you know the following:
\begin{itemize}
  \item your section number and your lecture schedules,
  \item your professors' office hours,
  \item the written assessments schedule,
  \item the online quizzes schedule,
  \item the midterm date and time.
\end{itemize}

\section*{Course Learning Outcomes}

Applied Mathematics 1201, Calculus \& Probability with Biological Applications, discuss mathematical modelling of natural phenomenons and analytic techniques of mathematical models with examples mostly chosen from biology.  The main classes of mathematical objects are
\begin{enumerate}
  \item power models \(y = c x^{\alpha}\),
  \item recurrence,
  \item ordinary differential equations (ODE),
  \item systems of recurrences and ODEs,
  \item probability, and discrete and continuous random variables,
  \item matricies, eigenvectors and eigenvalues (with complex numbers).
\end{enumerate}

Upon the completion of the course, you will be able to do the followings:

\begin{itemize}
  \item Describe natural phenomenons using power, discrete-time and continuous-time models.
  \item Describe uncertainty using probability and random variables.
  \item Discern appropriate applications of discrete-time versus continuous-time models.
  \item Discrete-time models.
  \begin{itemize}
    \item Describe natural phenomenon using recurrences.
    \item Solve linear recurrences.
    \item Find parameters in a linear recurrence satisfying a given condition.
    \item Describe the long-term behaviour of systems of linear recurrences.
  \end{itemize}
  \item Continuous-time models.
  \begin{itemize}
    \item Understand the exponential-growth, logistic models and the Lokta-Volterra models.
    \item Use ordinary differential equations to model mixing problems.
    \item Using
  \end{itemize}
  \item Solve discrete-time models up to three variables.
\end{itemize}

\section*{How to succeed in this courses}

\begin{enumerate}
  \item Study with purposes.
  
  Learning objectives describe expectations.  Every section in this textbook begins with a list of learning objectives.  Lecture and exercises are designed to help you meet learning objectives.  Be mindful of the learning objectives as you do exercises helps keep the targets in sight.
  \item Stroll to finish a marathon.
  
  A fourteen-week course is a marathon.  Lectures, exercises and formative assessments (quizzes and written assignments) are designed to help you develop \textbf{structured reasoning skills}, one stride at a time.  Digest concepts and develop skills with \textbf{consistent and evenly spaced out reviews and practice} sessions.
  \item Give yourself time to grow.
  \item Remember the precise statement of each definition and theorems.
  
  Definitions and theorem are precise statements that describe
\end{enumerate}

\end{document}
